\documentclass[border=4pt]{standalone}
\usepackage{amsmath,amssymb,mathtools,xcolor,varwidth}
\begin{document}
\begin{varwidth}{28em}
\large\bfseries
Draw the tangent $AD$, the straight line that just grazes the curve at $A$, matching its direction there and extending both ways. Between the tangent and the chord opens the angle $BAD$. As $B$ approaches $A$, the chord swivels toward the tangent, and this angle is diminished \emph{in infinitum}. For if the angle did \emph{not} vanish, the arc would meet the tangent at a finite rectilinear angle — a sharp elbow at $A$ — and the curvature there would not be continuous, against the supposition. Hence ultimately $\angle BAD \to 0$. \textit{Q.E.D.}
\end{varwidth}
\end{document}
